\documentclass[11pt]{article}
\usepackage{amsmath,hyperref}
\title{Currently unpaid does not mean the primary is terminal}
\date{}
\begin{document}\maketitle

\paragraph{Scope.}
An authorized consumer receives one fallback opportunity at time 0.4.
An earlier accepted primary can still arrive or apply after that cut.
This is an authored use of the byte-pinned source harness with a
separate controlled application queue. It does not claim that the
native network implementation delays its own immediate callback.
The first old report and independent work actually arrive in every
case, and the checked obligation here requires one actual application
settlement by 0.7, not one accepted command or one native transmit.

\paragraph{State and authority.}
A primary scheduled at 0.3, 0.4 or 0.6 can still be unattempted
at the alternative opportunity. At a tick-0.4 tie, order is explicit.
A successful physical primary receive can alternatively wait in the
authored application until 0.65 and then pay or fail; this is not a
remote guess from missing receipt. Primary TX costs three bytes and
fallback two, after an actual prior two-byte send. A failed physical
primary attempt may retain the three-byte charge. Only the specified
controller can revoke a queued attempt before it fires. A separately
authorized application owner can suppress a buffered application
effect, but cannot refund the native debit. A shared actor cannot
assume both permissions merely because it observes a clear status.
Fallback path and contact are fixed fixture capabilities.

A controlled unsafe comparison tries fallback whenever currently
unpaid. If the queued primary later pays and both attempts fit,
two actual settlements result. Waiting for actual terminal state
preserves safety but misses this opportunity when the primary has
not yet happened and will fail. If the source owns the still-queued
old operation, atomically checking fallback headroom, revoking that
old operation and taking the alternative obtains exactly one effect.
If primary's old event goes first at the same tick, the revoke is
no longer a source-queue cancellation. Suppressing a buffered
application would require the separate authority. Even if that
authority exists, an already charged transmit can exclude fallback.
The controller checks fallback ability before dropping the old option.

\paragraph{Method and limits.}
Actual native receive, queued downstream application and permitted
revocation remain distinct. The delayed application is executed
at the authored event time. Unrelated source callbacks are separately
checked. The direct same-information event checker agrees across
the enumerated native rows; there is no claim of a breakthrough at
matched resources. Not every observed unpaid state is distinguishable
from a late effect by the remote interpreter. This is a specialized
instance of the application boundary discussed in the primary
\href{https://pages.cs.wisc.edu/~pb/740/end-to-end.pdf}
{end-to-end arguments} and
\href{https://birrell.org/andrew/papers/ImplementingRPC.pdf}
{RPC implementation}. An independent test, actual authority
qualification and a protocol across uncoordinated owners are
still owed. Lean checks the small exactly-once and debit arguments,
not the C++ simulator or any deployed source.
\end{document}
